% Methodology
% Be specific enough that someone could reproduce your work

\label{sec:methodology}

% Overview
This study employs a mixed-method computational approach that bridges our 
\textit{motivating research questions}---concerning bias, neutrality over time, and democratic 
deliberation on Wikipedia---with the \textit{operationalized questions} that can be answered 
using observable Wikipedia data. Because constructs such as ``bias,'' ``deliberation,'' and 
``democratic participation'' cannot be measured directly, our methodology centers on identifying 
appropriate \textit{proxies} and applying machine learning tools to compare these signals across 
multiple Wikipedia Talk pages.

All data-collection code, processing scripts, and instructions for reproducing every figure in 
this paper are available in our GitHub repository:  
\href{https://github.com/brookeengland/Wikipedia-Bias-and-Democracy}{github.com/brookeengland/Wikipedia-Bias-and-Democracy}.

\subsection{Operationalizing Research Questions Through Proxies}

Our motivating questions ask how Wikipedia handles bias, how discussions evolve, and whether 
Talk pages reflect democratic deliberation. To answer these using measurable data, we identify 
the following proxies:

\begin{itemize}
    \item \textbf{Bias in content}: sentiment cues, lexical framing, and toxicity patterns  
    (following Hube's sentence-level bias indicators).
    \item \textbf{Deliberative quality}: proportion of Neutral vs.\ Positive/Negative comments.
    \item \textbf{Adversarial or harmful discourse}: toxicity levels (Low/Medium/High).
    \item \textbf{Democratic participation}: number of unique contributors and balance of participation 
    across editors.
\end{itemize}

These proxies are commonly used in prior research on Wikipedia governance and online 
deliberation, and they provide computable indicators aligned with our research questions.

\subsection{Data Collection}

Data were collected exclusively from Wikipedia Talk pages through the Wikipedia API using the 
\texttt{action=query} and \texttt{prop=revisions} endpoints. For each topic, we retrieved the 
full Talk page content and extracted individual comments using a timestamp-based regular 
expression matching the standard Wikipedia signature format.

Our updated analysis script supports the comparison of \textbf{multiple Talk pages at once}, 
allowing cross-topic evaluation (e.g., \textit{United States}, \textit{Climate change}, 
\textit{Elon Musk}).

\subsection{Data Processing}

Each Talk page undergoes the following preprocessing steps:

\begin{enumerate}
    \item \textbf{Comment segmentation}: isolate user comments based on signature patterns.
    \item \textbf{Cleaning}: remove usernames, timestamps, templates, and markup.
    \item \textbf{Chunking}: split long comments into sections that fit within model token limits.
    \item \textbf{Normalization}: store sentiment and toxicity results per-topic to enable direct comparison.
\end{enumerate}

Only the textual content of comments was processed, as non-textual metadata does not contribute 
to the deliberative or tonal characteristics of discussions.

\subsection{Analysis Methods}

We analyze each extracted comment using two transformer-based machine learning models:

\begin{itemize}
    \item \textbf{Sentiment Analysis}: Using the TabularisAI multilingual model, comments are categorized 
    as Positive, Neutral, or Negative. Neutral sentiment is treated as a proxy for reasoned, 
    deliberative communication.
    \item \textbf{Toxicity Detection}: Using the Unitary Toxic-BERT model, each comment is assigned a 
    toxicity level (Low, Medium, High).
\end{itemize}

For each Talk page, we aggregate:

\begin{itemize}
    \item sentiment label distributions,
    \item sentiment confidence levels,
    \item toxicity levels,
    \item stacked sentiment percentages.
\end{itemize}

The updated script automatically generates comparative visualizations across all selected topics, 
allowing us to operationalize questions such as:  
(1) whether politically contentious topics show higher toxicity,  
(2) whether some pages exhibit more democratic participation than others, and  
(3) how the tone of discourse varies by subject matter.

All machine learning analysis and visualization code is included in our repository and can be 
run with a single Python script.

\subsection{Reproducibility}

To ensure reproducibility, the repository contains:

\begin{itemize}
    \item full API querying code,  
    \item preprocessing and comment-extraction scripts,  
    \item instructions for installing Python dependencies,  
    \item the full multi-topic analysis script used to generate the results in this paper.
\end{itemize}

Not all implementation details appear in the PDF; however, the GitHub 
repository provides everything needed to regenerate our results exactly.

\subsection{Ethical Considerations}

All analyzed data were collected from publicly available Wikipedia Talk pages in accordance with 
\href{https://foundation.wikimedia.org/wiki/Policy:Terms_of_Use}{Wikimedia’s Terms of Use}. No private 
or user-identifiable information was collected or stored. The study complies with institutional 
guidelines for research using publicly accessible online data.